\documentclass[fontsize=10pt]{article}
\usepackage[margin=0.70in]{geometry}
\usepackage{lipsum,mwe,abstract}
\usepackage[T1]{fontenc}
\usepackage[english]{babel}

\usepackage{fancyhdr}
\setlength\parindent{0pt}

\usepackage{amsmath,amsfonts,amsthm}
\usepackage{wrapfig}
\usepackage{graphicx}
\usepackage{float}
\usepackage{subcaption}
\usepackage{comment}
\usepackage{enumitem}
\usepackage{cuted}
\usepackage{sectsty}
\allsectionsfont{\normalfont \normalsize \scshape}

\renewenvironment{abstract}
 {\small
  \begin{center}
  \bfseries \abstractname\vspace{-.5em}\vspace{0pt}
  \end{center}
  \list{}{%
    \setlength{\leftmargin}{0mm}
    \setlength{\rightmargin}{\leftmargin}%
  }
  \item\relax}
 {\endlist}

\makeatletter
\renewcommand{\maketitle}{\bgroup\setlength{\parindent}{0pt}
\begin{flushleft}
  \textbf{\@title}
  \@author \\
  \@date
\end{flushleft}\egroup
}
\makeatother

%% -------------------------------------------------------------------

\title{
\Large NS-RAM Co-Design: A Differentiable Software Bridge from Silicon to Algorithm \\
[10pt]
}
\date{\today}
\author{E. Bergvall, R. Luciani, S. Pazos, M. Lanza}

\begin{document}

 \maketitle

% --------------- ABSTRACT
\begin{abstract}
NS-RAM is a floating-body 2T cell whose intrinsic dynamics --- impact ionisation, parasitic
bipolar action, capacitive body retention --- can be coerced into acting as either a synapse,
a neuron, or a short-term memory element on demand. This proposal funds a closed-loop
co-design pipeline that lets the hardware team (Pazos, Lanza) test physical design choices
at the speed of software, and lets the algorithm team (Bergvall, Luciani) calibrate
ML primitives directly against silicon. The deliverable is a differentiable, fab-aware
device-level simulator that automatically produces a tape-out-ready design rule sheet for
the next TSMC mask, plus a headline demonstration of \emph{meta-plasticity} --- a single
NS-RAM device dynamically reassigning its computational role under control voltages.
\end{abstract}

\rule{\linewidth}{0.5pt}

% --------------- MAIN CONTENT

\section{Background}
NS-RAM cells exhibit a rich set of analog behaviours --- snapback, hysteretic body charge,
voltage-dependent capacitance --- that are intrinsic to the silicon and that conventional
digital memories actively suppress. Pazos and Lanza have shown that, with the right control
voltages, the same physical cell can be operated as a multi-bit weight, a leaky integrator,
or a stochastic spike generator. This makes NS-RAM a compelling substrate for in-memory
neuromorphic compute, but it also makes algorithm--device co-design unusually expensive:
each new circuit-level choice has to be verified against an opaque chain of SPICE
simulations, fabrication, and characterisation, with months of latency.

The bottleneck is not the silicon. It is the software bridge between the device physics
and the algorithm specification. Today, the algorithm team has no way to ask
``what control voltages would make this cell behave like a synapse with these statistics?''
without round-tripping through human-curated SPICE decks. Conversely, the hardware team has
no way to know whether a proposed transistor tweak will help or hurt a downstream learning
task without committing to a fabrication cycle.

\section{Objectives}
This proposal funds the construction of that bridge. Four concrete deliverables, in order of
risk:

\begin{enumerate}[leftmargin=1.5em,topsep=0.2em,itemsep=0.2em]
\item \textbf{Differentiable, fab-aware device simulator (Bergvall).} A PyTorch port of the
NS-RAM 2T cell that matches Pazos's measured DC and pulse responses to within publication
tolerance, exposes every model-card parameter as a leaf tensor with autograd enabled, and
can therefore be inverted: given a target algorithmic primitive (a bounded weight, a
specified time constant), the simulator returns the silicon parameters required to realise
it.
\item \textbf{Real-time topology simulator (Luciani).} An MTK-based engine that runs full
NS-RAM networks at $\geq 10^6$ cell-evaluations per second on a single workstation, allowing
multi-topology, multi-task sweeps in hours rather than weeks. Closes the loop with the
hardware team: a ``what-if'' question on a transistor parameter returns a
network-level performance number on the same day.
\item \textbf{Tape-out design rule sheet.} A short engineering document, derived from
sensitivity analysis on the differentiable simulator, that tells Pazos and Lanza which
physical tolerances matter for the target algorithms and which can be relaxed. De-risks
the next TSMC mask cycle and minimises silicon iterations.
\item \textbf{Headline demonstration: meta-plasticity.} A single NS-RAM device, operated
on a workstation through the differentiable simulator, dynamically reassigns its
computational role between synapse, neuron, and short-term memory under purely
electrical control --- end-to-end, in software, with the silicon and the
algorithm using identical parameters. This is the demo Lanza requires for the NRF
mid-sized grant submission.
\end{enumerate}

\section{Methodology}
The differentiable simulator is built on a faithful BSIM4 v4.8.3 implementation in
PyTorch, extended with the parasitic Gummel--Poon NPN, the deep N-well to body junction,
and a pseudo-arclength continuation solver that is robust to the snapback fold characteristic
of floating-body cells. Pazos's per-bias parameter table is honoured as the polynomial form
that ties model parameters to control voltages, so every sweep we run is parameterised on
the same axes the hardware team designs against. Luciani's MTK simulator imports the
calibrated cell as a callable subroutine; this guarantees that fast network-level sweeps
and slow, accurate device-level sweeps stay numerically consistent.

The two simulators are validated against three independent reference points: ngspice on
isolated transistors, Pazos's measured DC families and high-speed transients, and selected
LTSpice golden-reference decks supplied by the hardware team. The aim is not to replace
SPICE but to expose silicon physics as a differentiable, batched function suitable for
co-design loops with autograd.

\section{Status (May 2026)}
The DC fidelity milestone is closed. Our PyTorch BSIM4 port reproduces Pazos's 33-curve
measured I-V family with full Newton convergence at every sampled $(V_{G1}, V_{G2})$ bias.
Median residual against measurement is $1.00$~decade — matching the historical refit
baseline exactly, but reached purely through correct physics rather than ad-hoc
calibration. On isolated M2, the subthreshold $I_d$ ratio against ngspice is now
$1.05\times$ across $4$ decades of $V_{gs}$; the threshold-voltage gap is
$\le 1.5$~mV at $V_{ds} \in \{0.05, 0.5, 2.0\}$~V; DIBL matches within $0.3$~mV/V.

The $\sim 200\times$ tightening (from an initial $10.8\times$ subthreshold ratio
and $-57$~mV Vth gap) came from a bisecting-comparison campaign in which we built a
rebuilt-with-debug-printf instance of ngspice-42 alongside a literal C-to-Python
port of the BSIM4 device equations. The two acted as ground-truth oracles for our
production pyport. Five separate calibration-loop bugs surfaced and were
documented:

\begin{enumerate}\setlength{\itemsep}{1pt}
  \item ngspice silently drops the 2nd and 3rd \texttt{key=value} assignments on
        any \texttt{.model}-body line, losing \texttt{wvth0}, \texttt{voff}-binning,
        \texttt{ww}, \texttt{binunit=2}, and similar parameters whenever the card
        author packs them with leading entries.
  \item \texttt{toxe = toxn} \texttt{.param} substitutions in the \texttt{.model}
        body silently fall back to the BSIM4 default (3~nm), not to the
        \texttt{.param} value (4~nm).
  \item \texttt{lpe0 = lpe0n} fails the same way, defaulting to 1.74e-7~m
        instead of the card's 1.244e-7~m; this term enters the lateral-pocket
        $V_{th}$ correction and accounts for $\sim 8$~mV of Vth offset in our cell.
  \item Our prior $\phi$ surface-potential formula carried an extra factor of two
        (a misreading of the BSIM4 manual); the correct ngspice form is
        $\phi = V_{tm} \log(N_{dep}/n_i) + \phi_{in} + 0.4$.
  \item \texttt{phin} is silently dropped despite being first on its card line,
        making the $+0.4$ constant in the $\phi$ formula the only contributor in
        cards that try to set \texttt{phin}.
\end{enumerate}

Pyport now mimics ngspice-42's calibration-loop behavior \emph{exactly}, while
preserving correct multi-assignment parsing for newly-authored cards. The bug
catalogue is publishable as a standalone note and, more practically, lets any
future card-fitting workflow against ngspice be reproduced in autograd.

The transient milestone is unblocked. Pazos's $5$--$10$~fF body-junction diode is
integrated and the implicit-Euler joint Newton solver converges with line search and
hard bounds at $100\%$ of randomly sampled $(V_{G1}, V_{G2})$ pairs. Each cell runs at
$\sim 14$~s per $500$-step trajectory on CPU; the body-cap time constant is
$\tau_{\mathrm{body}} \approx 0.7$~ns, which sets the natural reservoir bandwidth.

The MTK topology simulator hits $5 \times 10^6$ cell-evaluations per second on a
single AMD Radeon~8060S iGPU using PyTorch ROCm with \texttt{torch.compile} on the
inner Newton kernel ($N=10^5$ cells, $40$-point Vd~sweep, $0.10$~s wall). This is
$5\times$ Mario's stated G2 throughput target. The transient driver remains CPU-bound
because of sequential time-stepping; per-cell wall is dominated by Python overhead
rather than physics, so a multi-cell threading layer brings the network-scale brief
into the same workstation budget without GPU porting of the transient solver.

\paragraph{First B.5 result --- memory-capacity validation under external recurrence.}
With every cell isolated, the linear-memory-capacity benchmark returns
$\mathrm{MC} = 0.22 \pm 0.10$ (mean $\pm$ std over 5 seeds, $N=10$ cells, $T=300$),
confirming the cell as a memoryless analog weight rather than a reservoir node.
Adding external CMOS-style recurrent routing
($V_{G2}^{\mathrm{eff}}[i] = V_{G2}^{0}[i] + \kappa \sum_j W^{\mathrm{rec}}_{ij}\,
\log_{10}|I_d^{(j)}[t-1]|$, $\kappa = 0.03$, echo-state init for $W^{\mathrm{rec}}$)
lifts $\mathrm{MC}$ to $\mathbf{1.10 \pm 0.23}$, a paired $\Delta = +0.88$ at
$t = 7.4$ over the same 5 seeds. This is the first quantitative demonstration on
the calibrated stack that the cell-as-weight + external-recurrence architectural
framing yields reservoir-quality temporal computation, and it sets the protocol
for the remaining four B.5 benchmarks (Hopfield, NARMA-10, temporal-XOR,
multi-class waveform), each to be reported with the same 5-seed paired-$t$ framing.

\paragraph{Second B.5 result --- temporal-XOR$(\tau=2)$.}
Under the same protocol with bipolar input $u(t) \in \{-1, +1\}$ and target
$y(t) = u(t) \oplus u(t-2)$, the linear sign-readout reaches accuracy
$0.54 \pm 0.04$ (5 seeds, chance $0.50$) without recurrence and
$\mathbf{0.68 \pm 0.13}$ with $\kappa=0.03$, a paired $\Delta = +0.13$ at
$t = +2.7$. Per-seed accuracies span $0.49$--$0.80$ at $\kappa=0.03$, a
characteristic echo-state-lottery distribution that the multi-seed report
captures honestly. NARMA-10 at the same $(N, T)$ scale is too hard for a
$10$-cell reservoir with a small input gain; we defer it to a $N{\geq}100$
rerun with a spectral-radius-controlled $W^{\mathrm{rec}}$ rather than report
a null at the wrong scale.

\paragraph{Third B.5 result --- Hopfield-style associative recall, with an informative asymmetry.}
With three stored $\pm 1$ prototype patterns and $20\%$ bit-flip corruption on
each cue, a one-vs-rest ridge readout over per-cell $\log_{10}|I_d|$ reaches
accuracy $\mathbf{0.69 \pm 0.07}$ (5 seeds, chance $0.33$) \emph{without} recurrence.
Adding $\kappa=0.03$ \emph{drops} accuracy to $0.58 \pm 0.09$ (paired
$\Delta = -0.11$, $t = -2.45$). The asymmetry between this benchmark and the
temporal ones above is the most architecturally interesting finding so far:
the cell array's per-cell nonlinearity already provides enough separation to
discriminate spatial patterns at modest noise, and cross-cell recurrence then
\emph{removes} the decoupled per-cell channels the readout relies on. External
recurrence is therefore \emph{task-class dependent} --- a benefit for tasks
demanding temporal memory, a liability for tasks rewarding decoupled per-cell
nonlinear channels. The architectural recommendation is task-conditioned routing
rather than blanket recurrence; this directly informs the C.3 tape-out
recommendation we will deliver to Mario for the next mask cycle.

\paragraph{Fourth B.5 result --- NARMA-10 at $N=100$, $\rho=0.9$.}
With $W^{\mathrm{rec}}$ spectral-radius-controlled to $\rho=0.9$ and a finer
$\kappa$ sweep, the stable operating point is $\kappa=0.003$: NRMSE drops
from $1.073 \pm 0.048$ at $\kappa=0$ to $\mathbf{0.946 \pm 0.018}$ (paired
$\Delta = -0.128$, $t = -9.4$, $T=600$, 5 seeds) --- the largest effect size
and tightest std across the four B.5 benchmarks. The chaos-onset boundary
sits between $\kappa=0.003$ and $\kappa=0.005$: one in five seeds diverges at
$\kappa=0.005$. The silicon shared-bulk-rail topology in C.3 therefore needs
a digitally tunable coupling resistor so the operating point can be scanned
in measurement.

\section{Limitations and open questions}
We summarise the visible caveats so reviewers see the honest scope of what is
and is not yet established:
\begin{itemize}\setlength{\itemsep}{1pt}
  \item \textbf{NARMA-10 is solved at the right operating point but is
    $\kappa$-sensitive.} At $\kappa=0.003$, $N=100$, $\rho=0.9$ we measure
    NRMSE $0.946 \pm 0.018$ (paired $t=-9.4$ vs $\kappa=0$); at $\kappa=0.005$
    one in five seeds diverges (NRMSE $2.21$); at $\kappa=0.03$ all seeds
    over-drive. This $\kappa$-bracket is narrow, so the silicon implementation
    must allow scanning. The $N=10$ regime remains too small for NARMA-10
    regardless of $\kappa$.
  \item \textbf{Hopfield reported at $N=10$, $M=3$ patterns only.}
    "Associative memory at scale" requires $N{\geq}50$, $M{\geq}30$; the
    $N$-scaling at $\kappa=0$ to confirm the substrate-alone advantage holds
    is on the queue and explicitly named as risk \#2 in the C.3 tape-out
    recommendation document.
  \item \textbf{Software vs.\ silicon recurrence equivalence is unproven.}
    The $\kappa$-mediated MC lift (0.22~$\to$~1.10) was demonstrated via
    \emph{external} $W^{\mathrm{rec}}$ in software. Whether an on-die
    shared-bulk-rail topology reproduces the same lift in silicon is the
    central scientific risk; the M9 fan-out test structure is exactly that
    experiment.
  \item \textbf{Thick-oxide cell card is pending from Pazos (A.12).}
    All thick-ox claims in this brief lean on PTM extrapolation rather than
    measured data. The M3 milestone re-fit is contingent on the data drop.
  \item \textbf{The five ngspice-42 calibration-loop bugs are documented but
    not yet upstreamed.} A short standalone note + a reproduction harness
    (debug-printf-instrumented build + literal C-to-Python port) is the
    natural next disclosure step; we will share it with the SimuCAD/ngspice
    maintainers once the quadrant of fixes (parser, $\phi$ formula,
    \texttt{lpe0}/\texttt{toxe} substitution, \texttt{phin}) is reviewed by
    Pazos.
\end{itemize}

\section{Deliverables and Timeline}
\textbf{M3 (Jun 2026)}: closure of the residual $\sim 10$~mV Vth gap; refit of
the M1 card against the latest measurement set. Acceptance: $<0.5$~decade median
on the 33-bias regression.\\
\textbf{M6 (Sep 2026)}: full B.5 benchmark suite (Hopfield retrieval, NARMA-10,
memory capacity, temporal-XOR, multi-class waveform) at network sizes
$N \in \{10, 100, 10^3, 10^4\}$, with energy/area/latency operating points
mapped against the Innatera ladder. Acceptance: at least one benchmark within
$1.5\times$ of the Innatera $10$--$100$~mW gateway baseline.\\
\textbf{M9 (Dec 2026)}: fan-out / shared-rail validation circuit
($10$--$30$ cells, per Pazos's request), feeding a tape-out-ready cell-parameter
recommendation for the next NS-RAM prototype.\\
\textbf{M12 (Mar 2027)}: thick-oxide $17$~$\mu$m$^2$ neuron-cell variant
(VG2~$\in [2.5, 3.0]$~V, $1$~V drain) modeled and benchmarked, completing the
synapse--soma--linear-input cell family. Final write-up.

\section{Budget}
Compute: existing Daedalus and Ikaros workstations are sufficient; no cloud
dependency. Personnel: $0.4$ FTE software (Bergvall) for the topology layer
and benchmarks, $0.2$ FTE consult time (Pazos, Lanza) for measurement
turnaround, optional $0.3$ FTE collaborator (Luciani, Julia cross-validation)
for the joint-newton kernel re-implementation in Julia/Enzyme, the latter
unlocking native AD on top of the existing PyTorch port. Travel: one
in-person tape-out review at KAUST in Q3 2026.

\begin{figure}[h]
  \centering
  \includegraphics[width=0.92\textwidth]{../figures/quadrant_nsram_vs_edge.pdf}
  \caption{NS-RAM 2T cell positioning against published edge-AI accelerator
  baselines. Energy axis = published MAC/synop energy; latency axis = vendor
  inference time for KWS-class workloads. NS-RAM sits in the lower-left
  ``ultra-low energy, ultra-low latency'' quadrant, three to five orders below
  digital edge processors. The dashed line marks the Innatera $10$--$100$~mW
  gateway envelope which Mario's NRF brief targets.}
  \label{fig:quadrant}
\end{figure}

\section{Conclusion}
The NS-RAM 2T floating-body cell sits at an energy point ($21$~fJ/cycle,
$6.7$~fJ/spike, $46$~$\mu$m$^2$) two orders of magnitude below digital
neuromorphic baselines while remaining standard-CMOS compatible. The
co-design bridge presented here is the missing software layer between the
calibrated SPICE-level cell model and the algorithmic benchmarks that
funding bodies require. Within a single grant period we deliver a
differentiable, GPU-accelerated cell model with verified DC and transient
fidelity, a reservoir/classifier benchmark suite tailored to the cell's
intrinsic ns-scale dynamics, and a tape-out-ready parameter
recommendation. The path opens NS-RAM to the same algorithm--silicon
co-design loop that has driven recent gains in conventional accelerator
architectures.


\begin{thebibliography}{}
\bibitem{Ref1}{Reference 1 -- Pazos et al., NS-RAM cell measurements}
\bibitem{Ref2}{Reference 2 -- Lanza et al., 2T floating-body memristor}
\bibitem{Ref3}{Reference 3 -- BSIM4 v4.8.3 manual, Berkeley}
\end{thebibliography}


 \end{document}
